% =========================================================
% Final integrative project -- Time Series Analysis module
% Author : Aymen Ben Brik (aymen.benbrik@esprit.tn)
% Esprit School of Business
% =========================================================
\documentclass[a4paper,11pt]{article}

\newcommand{\TPnumero}{Final integrative project}
\newcommand{\macrosPath}{../preamble/macros.tex}
% =========================================================
% Tutorial / lab preamble — Time Series Analysis module
% article class + fancyhdr + tcolorbox + listings (English)
% Author : Aymen Ben Brik (aymen.benbrik@esprit.tn)
% =========================================================

\usepackage[utf8]{inputenc}
\usepackage[T1]{fontenc}
\usepackage[english]{babel}
\usepackage{lmodern}
\usepackage{amsmath, amssymb, amsthm, mathtools, bm}
\usepackage{graphicx}
\usepackage{booktabs}
\usepackage{array}
\usepackage{multirow}
\usepackage{colortbl}
\usepackage{xcolor}
\usepackage{tikz}
\usetikzlibrary{positioning, arrows.meta, calc, shapes, fit, backgrounds, matrix}
\usepackage{listings}
\usepackage[most]{tcolorbox}
\usepackage{geometry}
\geometry{a4paper, margin=2cm, headheight=15pt}
\usepackage{fancyhdr}
\usepackage[hidelinks]{hyperref}
\usepackage{enumitem}

% --- Colours ---
\definecolor{espritBlue}{RGB}{30, 60, 110}
\definecolor{espritAccent}{RGB}{200, 80, 30}
\definecolor{boxEx}{RGB}{30, 60, 110}
\definecolor{boxQ}{RGB}{180, 120, 0}
\definecolor{boxC}{RGB}{30, 100, 60}
\definecolor{boxAtt}{RGB}{200, 80, 30}

% --- Listings ---
\definecolor{codebg}{RGB}{248,248,248}
\definecolor{codeframe}{RGB}{220,220,220}
\definecolor{keyword}{RGB}{0,82,155}
\definecolor{comment}{RGB}{88,110,117}
\definecolor{string}{RGB}{133,153,0}

\lstdefinestyle{pythonstyle}{
  language=Python,
  basicstyle=\ttfamily\small,
  keywordstyle=\color{keyword}\bfseries,
  commentstyle=\color{comment}\itshape,
  stringstyle=\color{string},
  backgroundcolor=\color{codebg},
  frame=single,
  rulecolor=\color{codeframe},
  frameround=tttt,
  breaklines=true,
  showstringspaces=false,
  tabsize=2
}
\lstdefinestyle{Rstyle}{
  language=R,
  basicstyle=\ttfamily\small,
  keywordstyle=\color{keyword}\bfseries,
  commentstyle=\color{comment}\itshape,
  stringstyle=\color{string},
  backgroundcolor=\color{codebg},
  frame=single,
  rulecolor=\color{codeframe},
  frameround=tttt,
  breaklines=true,
  showstringspaces=false,
  tabsize=2
}
\lstset{style=pythonstyle}

% --- Common macros (configurable path) ---
\providecommand{\macrosPath}{../../preamble/macros.tex}
\input{\macrosPath}

% --- Exercise counter ---
\newcounter{excounter}
\renewcommand{\theexcounter}{\arabic{excounter}}

\newtcolorbox[use counter=excounter]{exercise}[1][]{
  enhanced, breakable,
  colback=boxEx!5, colframe=boxEx, fonttitle=\bfseries,
  title={Exercise~\theexcounter\ifstrempty{#1}{}{ -- #1}},
  arc=2pt, boxrule=0.6pt
}
\newtcolorbox{question}[1][]{
  enhanced, breakable,
  colback=boxQ!5, colframe=boxQ, fonttitle=\bfseries,
  title={Question\ifstrempty{#1}{}{ -- #1}},
  arc=2pt, boxrule=0.5pt
}
\newtcolorbox{solution}[1][]{
  enhanced, breakable,
  colback=boxC!5, colframe=boxC, fonttitle=\bfseries,
  title={Solution\ifstrempty{#1}{}{ -- #1}},
  arc=2pt, boxrule=0.5pt
}
\newtcolorbox{warning}[1][]{
  enhanced, breakable,
  colback=boxAtt!5, colframe=boxAtt, fonttitle=\bfseries,
  title={Warning\ifstrempty{#1}{}{ -- #1}},
  arc=2pt, boxrule=0.5pt
}

% --- Headers / footers ---
\pagestyle{fancy}
\fancyhf{}
\fancyhead[L]{\textsc{\moduleNom} -- \auteurNom}
\fancyhead[R]{\TPnumero}
\fancyfoot[C]{\thepage}
\fancyfoot[L]{\auteurInstitution}
\fancyfoot[R]{\auteurEmail}
\renewcommand{\headrulewidth}{0.4pt}
\renewcommand{\footrulewidth}{0.2pt}

% --- Placeholder ---
\providecommand{\TPnumero}{Lab}


\title{Final integrative project\\[2pt]
\large End-to-end time-series pipeline\\
\large Statistics $\cdot$ Decomposition $\cdot$ ARMA $\cdot$ GARCH $\cdot$ VaR}
\author{\auteurNom\\\auteurEmail\\\auteurInstitution}
\date{\anneeUniversitaire}

\begin{document}
\maketitle

% =========================================================
\section*{Context and objectives}

This integrative project mobilises the notions of every chapter of
the module:
\begin{itemize}
  \item \textbf{Ch.~1 Basic statistics}: descriptive analysis,
        confidence intervals, normality tests.
  \item \textbf{Ch.~2 Decomposition}: classical time-series
        decomposition.
  \item \textbf{Ch.~3 Stationarity}: ACF/PACF and Ljung--Box.
  \item \textbf{Ch.~4 Non-stationarity}: ADF/KPSS unit-root tests
        and differencing.
  \item \textbf{Ch.~5 ARMA}: Box--Jenkins identification, MLE,
        forecasting.
  \item \textbf{Ch.~6 ARCH/GARCH}: Engle LM, GARCH$(1, 1)$, and
        Value-at-Risk.
\end{itemize}

\noindent
The fil-rouge task is a complete \emph{risk-management study} of a
synthetic but realistic daily stock-price series. The goal is to
produce, by the end of the project, a one-day Value-at-Risk that is
well-calibrated according to a Kupiec back-test.

\section*{Setup}
\begin{itemize}
  \item Language: Python 3, dependencies: \texttt{numpy},
        \texttt{pandas}, \texttt{matplotlib}, \texttt{scipy},
        \texttt{statsmodels}, \texttt{arch}.
  \item Reproducibility: \texttt{numpy.random.seed(42)}.
  \item Skeleton in \texttt{seed/projet.py}; full reference
        solution in \texttt{correction/projet.py}.
  \item Dataset: \texttt{data/stock\_prices.csv} -- $T = 2000$
        synthetic daily prices with drift and GARCH volatility,
        regenerated by \texttt{generate\_dataset.py}.
  \item Deliverables:
    \begin{itemize}
      \item a reproducible \texttt{projet\_<name>.py};
      \item a 10--12 page report PDF with figures and discussion;
      \item a 12-min individual oral defence.
    \end{itemize}
\end{itemize}

% =========================================================
\section*{Data}

The data file \texttt{data/stock\_prices.csv} contains $T = 2000$
daily closing prices $P_t$ over the period 2018-01-02 to 2025-09-15
(business days). The series is generated as
\[
  \log P_t = \log P_{t-1} + \mu + \varepsilon_t,
  \qquad
  \varepsilon_t = \sigma_t z_t,
  \quad z_t \iid \mathcal N(0, 1),
\]
with $\sigma_t^2 = \alpha_0 + \alpha_1 \varepsilon_{t-1}^2 + \beta_1
\sigma_{t-1}^2$ -- a textbook GARCH$(1, 1)$ data-generating process
on the log-returns. The true parameters are kept hidden; you will
recover them from the data.

% =========================================================
\begin{exercise}[Part A -- Basic statistics on returns (Ch.~1)]
\textbf{(A.1)} Load \texttt{data/stock\_prices.csv}. Compute the
log-returns $r_t = \log(P_t / P_{t-1})$, $t = 1, \dots, T-1$.

\textbf{(A.2)} Report the sample mean $\widehat\mu$, sample
standard deviation $\widehat\sigma$, sample skewness, sample
excess kurtosis. Provide a $95\%$ confidence interval for the
mean (using the Student-$t$ approximation).

\textbf{(A.3)} Test the null hypothesis $H_0: \mu = 0$ at the $5\%$
level. Conclude.

\textbf{(A.4)} Plot a histogram of $r_t$ together with the
$\mathcal N(\widehat\mu, \widehat\sigma^2)$ density. Run a
Jarque--Bera normality test and comment.
\end{exercise}

% =========================================================
\begin{exercise}[Part B -- Decomposition of the price level (Ch.~2)]
\textbf{(B.1)} On the price series $P_t$ (\emph{not} the returns),
extract the trend by a centred moving average of width $k = 50$
business days.

\textbf{(B.2)} On the log-price $\log P_t$, fit a linear trend by
OLS:
\[
  \log P_t = \beta_0 + \beta_1 t + u_t.
\]
Report $\widehat\beta_1$ -- the implied mean log-return. Compare
with the sample mean of $r_t$ (Part A).

\textbf{(B.3)} Plot $P_t$, the moving-average trend and the OLS
exponential trend on the same figure. Comment in 5 lines.
\end{exercise}

% =========================================================
\begin{exercise}[Part C -- Stationarity diagnostics (Ch.~3)]
\textbf{(C.1)} Plot the sample ACF of $\log P_t$ up to lag~$60$.
What do you observe? Is the series consistent with a stationary
process?

\textbf{(C.2)} Same on the returns $r_t$. Are returns serially
uncorrelated?

\textbf{(C.3)} Run the Ljung--Box test on $r_t$ at lags $m \in
\{10, 20\}$. Conclude.

\textbf{(C.4)} Run the Ljung--Box test on $r_t^2$ at the same
lags. What do the two results, taken together, suggest?
\end{exercise}

% =========================================================
\begin{exercise}[Part D -- Unit-root testing (Ch.~4)]
\textbf{(D.1)} Run the ADF test on $\log P_t$ with
\texttt{regression="ct"} (constant + trend) and on $r_t$ with
\texttt{regression="c"}. Report stat, $p$-value, decision at
$5\%$ in a small table.

\textbf{(D.2)} Run the KPSS test on the same two series.

\textbf{(D.3)} Build the ADF/KPSS combination matrix on
$\log P_t$. In which cell does it land? What does this imply
about the integration order $I(d)$ of the price series?

\textbf{(D.4)} Confirm: $\log P_t \sim I(1)$ and $r_t \sim I(0)$.
\end{exercise}

% =========================================================
\begin{exercise}[Part E -- ARMA modelling (Ch.~5)]
\textbf{(E.1)} On the returns $r_t$, plot the sample ACF and PACF
up to lag~$20$. Identify any structure in the conditional mean.

\textbf{(E.2)} Fit three candidates by MLE
(\texttt{statsmodels.tsa.arima.ARIMA}):
\begin{itemize}
  \item ARMA$(0, 0)$: constant mean only;
  \item ARMA$(1, 0)$;
  \item ARMA$(0, 1)$.
\end{itemize}
Report a comparison table: log-likelihood, AIC, BIC,
Ljung--Box($24$) $p$-value.

\textbf{(E.3)} Choose the best model on the (AIC, BIC, Ljung--Box)
triplet. The constant-mean model often wins -- justify why this is
expected for daily returns.
\end{exercise}

% =========================================================
\begin{exercise}[Part F -- GARCH$(1, 1)$ and VaR (Ch.~6)]
\textbf{(F.1)} On the residuals $\widehat\varepsilon_t$ of the
chosen mean model from Part~E, run Engle's LM test at lags $5, 10,
20$. Confirm the presence of ARCH effects.

\textbf{(F.2)} Fit a GARCH$(1, 1)$ on the returns $r_t$ with a
constant mean and Gaussian innovations using the
\texttt{arch} package. Report
$\widehat\mu$, $\widehat\alpha_0$, $\widehat\alpha_1$,
$\widehat\beta_1$ with their standard errors. Compute
$\widehat\alpha_1 + \widehat\beta_1$ and the implied
unconditional variance.

\textbf{(F.3)} Compute the standardised residuals $\widehat z_t =
\widehat\varepsilon_t / \widehat\sigma_t$. Run Engle's LM test on
$\widehat z_t^2$. Has the GARCH absorbed the volatility clustering?

\textbf{(F.4)} Compute the conditional one-day Value-at-Risk
\[
  \mathrm{VaR}_{T+1}(\alpha)
  = -\bigl(\widehat\mu + z_\alpha\, \widehat\sigma_{T+1}\bigr),
  \qquad
  \alpha \in \{1\%, 5\%\}.
\]
Compare with the unconditional VaR (using the sample~$\widehat\sigma$).
\end{exercise}

% =========================================================
\begin{exercise}[Part G -- Kupiec back-test (synthesis)]
\textbf{(G.1)} Split the data: train = first $1750$ obs, test =
last $250$ obs.

\textbf{(G.2)} On the test set, compute a \emph{rolling}
conditional VaR at $\alpha = 5\%$: at each day $t$ in the test set,
fit the GARCH on observations $1..t-1$, compute
$\widehat\sigma_t$, and predict $\mathrm{VaR}_t(5\%)$.

\textbf{(G.3)} Count the number of \emph{exceedances} -- days where
the realised loss $-r_t$ exceeded $\mathrm{VaR}_t(5\%)$.

\textbf{(G.4)} Run the Kupiec proportion-of-failures test:
\[
  LR_{\text{Kupiec}}
  = -2 \log \left[
      \frac{(1 - \alpha)^{N - X} \alpha^X}
           {(1 - \widehat\pi)^{N - X} \widehat\pi^X}
    \right]
  \;\xrightarrow{d}\; \chi^2_1,
\]
with $N = 250$, $X$ the number of exceedances, $\widehat\pi = X / N$.
Reject the null of correct unconditional coverage if
$LR > \chi^2_{1, 0.95} = 3.841$.

\textbf{(G.5)} Conclude: is the GARCH-VaR well-calibrated?
\end{exercise}

% =========================================================
\section*{Marking grid (out of 100)}

\begin{center}
\begin{tabular}{p{8cm}rl}
\toprule
\textbf{Criterion} & \textbf{Points} & \\
\midrule
Part A: descriptive statistics + Jarque--Bera     & 10 \\
Part B: decomposition + linear-trend fit          & 10 \\
Part C: ACF/PACF + Ljung--Box on $r_t$ and $r_t^2$ & 10 \\
Part D: ADF/KPSS + combination matrix             & 15 \\
Part E: ARMA candidates + AIC/BIC/LB              & 15 \\
Part F: GARCH$(1,1)$ + standardised diagnostics + VaR & 20 \\
Part G: Kupiec back-test                          & 10 \\
Reproducibility + report quality                  &  5 \\
Oral defence                                      &  5 \\
\midrule
\textbf{Total}                                    & \textbf{100} \\
\bottomrule
\end{tabular}
\end{center}

\medskip
\noindent\textbf{Bonus.} Student-$t$ innovation distribution in the
GARCH (\(+5\) points), Christoffersen's conditional-coverage test
on top of Kupiec ($+5$), comparison with EGARCH or GJR-GARCH
($+5$).

% =========================================================
\section*{Submission}

\begin{itemize}
  \item \textbf{Code}: \texttt{projet\_<name>.py} (or
        \texttt{projet\_<name>.ipynb}) reproducible with seed $42$.
  \item \textbf{Report}: 10--12 PDF pages including:
        \begin{itemize}
          \item summary of the dataset and stylised facts;
          \item all required tables and figures (one figure per
                part);
          \item Kupiec back-test result;
          \item discussion of model choice and calibration.
        \end{itemize}
  \item \textbf{Oral defence}: 12 minutes per student, 5 minutes
        of questions.
\end{itemize}

\bigskip
\begin{warning}
Install once: \texttt{pip install numpy scipy matplotlib pandas
statsmodels arch}. The full reference pipeline runs in
$\approx 15$ seconds on a laptop CPU; the rolling Kupiec
back-test (Part G) is the slowest part ($\sim 30$~s).
\end{warning}

\bigskip
\noindent\textit{Good luck!}\par\medskip
\hfill\auteurNom\par
\hfill\auteurEmail\par

\end{document}
